\scanpage{156}
\mspage{30}
% archive sheet number: 155

dont les \unsure{quatre} premières utilisent seulement
l'hypothèse $\Phi\longleftrightarrow\Psi$ (de a) de la définition
4.1.), et la dernière seulement, $d\Rightarrow c'$, fait appel
à l'hypothèse de factorisation de Quillen (b) de
4.1.) ; et encore il suffit de la supposer pour $j$ seulement !.

$c')\Rightarrow c)$ tautologie pure, $c)\Rightarrow a)$ car on sait
que $\widetilde\Phi$ est stable par facteurs directs (et contient $\Phi$),
$a\Rightarrow b$ car $\widetilde\Phi\subset\Psi^*$ ($\widetilde\Phi$ étant par définition la plus
petite partie close à g. de $\mathrm{Fl}\,M$ contenant $\Phi$,
et $\Psi^*$ étant close à gauche), $b\Rightarrow d$ car on écrit
le diagramme donné 4.2.2. (traits pleins)
comme

\diagram{\begin{tikzcd}[ampersand replacement=\&,column sep=large]
\& A \arrow[r,"\unsure{u}"] \arrow[d,"j"',"\in\Psi^*"] \& X \arrow[d,"p\in\Psi"]\\
\text{($\unsure{\alpha|a|\sigma}$)} \& B \arrow[r,"v"'] \arrow[ur,dashed,"q"'] \& Y
\end{tikzcd}}
en y faisant \deleted{\unsure{$B=X$}} $Y=B$, $v=\mathrm{id}_B$ \inserted{et en rebaptisant $X$ \deleted{\unsure{de}} par $B'$, et $u$ par $i$ :} ;
l'hypothèse b) $j\in\Psi^*$ signifie que dans un diagramme \inserted{traits pleins} ($\unsure{\alpha|a}$), il
existe tjrs $q$ telle que $qj=u$, $pq=v$,
et en l'exprimant dans le cas particulier
où $v=\mathrm{id}_B$, on trouve d).

Reste à prouver $d)\Rightarrow c')$, en utilisant
cette fois l'hypothèse de factorisation sur $j$, i.e.
qu'on peut écrire $j$ sous la forme $pi$ :
\diagram{\begin{tikzcd}[ampersand replacement=\&,column sep=small]
 A \arrow[rr,"{\unsure{\Delta|j}}"] \arrow[dr,"i"'] \& \& B \arrow[dl,bend left=65,dashed,"q"]\\
 \& B' \arrow[ur,"p"'] \& {,}
\end{tikzcd}}
avec $i\in\Phi$, $p\in\Psi$. En vertu de
d), \inserted{comme $p\in\Psi$} il existe donc $q:B\longrightarrow B'$ satisfaisant
